\abstractPT  % Do NOT modify this line

Este trabalho apresenta quatro solucionadores automáticos para o puzzle Sudoku clássico e para a sua variante Killer Sudoku. Em Prolog, foram implementados o algoritmo de aprofundamento iterativo (DFID) e o algoritmo A*, ambos com a heurística de valores mínimos restantes (MRV). Em Python, foram implementados o arrefecimento simulado (SA) e o algoritmo genético (GA), com representação baseada em trocas intra-caixa. Todos os solucionadores foram estendidos para suportar Killer Sudoku. Os resultados mostram que o DFID com MRV resolve todas as instâncias em menos de quatro segundos. O A* expande mais nós em puzzles com elevado grau de liberdade. O SA resolve de forma fiável instâncias fáceis e intermédias, ficando por vezes preso em mínimos locais de custo dois. O GA converge prematuramente e não encontra soluções completas dentro dos limites testados.
